\section{Recommendation and review method}
\label{sec:recommendation}
\paragraph{Recommended first paper.}
Ask: \emph{When does spatial Dysthe physics improve reconstruction and
downstream prediction from sparse water-wave observations, after NLS is
allowed to fit effective coefficients?} The contribution should be a
validated map of when the additional physics is useful, including
uncertainty and computational cost. It should not rest solely
on substituting a more complicated residual into a familiar neural network.

The clearest published motivation is Ehlers et al.'s water-wave NLS study,
whose Conclusions propose Dysthe as a next model \cite{ehlers2024}.
However, subsequent papers already address fully nonlinear wave
reconstruction and reusable physics-informed operators
\cite{ehlers2025,wavepino}. Therefore, generic claims of the first PINN
for water-wave reconstruction or a fast water-wave PINO are unsuitable.
The more defensible opportunity is a controlled comparison of physics,
identifiability, phase accuracy, and generalization in the specified equation.

\paragraph{What newest and best mean here.}
This is a targeted critical review of primary literature available by
25 October 2026, not an exhaustive systematic review or a universal ranking.
Searches combined NLS/Schr\"odinger, PINN, water waves, Dysthe,
conservation, causal training, optimization, error bounds, sparse
assimilation, and operator learning; recent searches emphasized 2024--2026.
Publisher records, author manuscripts, proceedings, author code, and data
repositories were consulted. Publication status was checked separately from
arXiv submission dates. The source ledgers supplied with the report record
access and limitations. A search failure is not proof of absence.

Priority is based on relevance to the task, quality of comparisons,
independent test data, access to methods and code, and clarity about costs.
There is no common benchmark establishing one best NLS--PINN method across
forward, inverse, long-distance, and many-query problems. Numerical errors
reported on different domains, norms, equations, and data budgets are not a
leaderboard. Very recent publisher abstracts are useful leads, but provide
less inspectable evidence than a complete manuscript and executable release.

\paragraph{Reading priorities.}
Read Ehlers et al. (2024) for the direct research motivation;
Urb\'an et al. (2025) for a serious compact NLS training baseline;
Zhang et al. (2025) for complex-valued NLS learning;
Hao et al. (2026) for hard-constraint ablations;
and the 2025 water-wave PINO paper for the many-query comparison
\cite{ehlers2024,urban,cvpinn,hardic,wavepino}.
Read the 2026 GRAVIS, Noether, and phase--amplitude papers as recent
method candidates, with the access limitations in Section~\ref{sec:literature}.
For a mathematics-led direction, add Alejo et al. and Eiras et al.
\cite{alejo,eiras}.

\paragraph{Evidence labels.}
Throughout this report, \emph{reported} means an author's result, not a result
reproduced here. \emph{Explicit future work} identifies an actual statement
in the inspected source. \emph{Proposed question} is this review's synthesis.
An available code repository is evidence of availability, not evidence that
the code reproduces every claim. The figures in this report are original
analytic illustrations and experiment diagrams, not new PINN experiments.

\paragraph{Scope.}
The literature survey concerns NLS methods broadly. All proposed active
research stays within unidirectional deep-water gravity waves and the
one-profile-coordinate spatial free-surface model
\texttt{water-wave-dysthe-spatial-v1}. General NLS, higher-order NLS, and
other Schr\"odinger demonstrations can motivate algorithms; they are not
water-wave Dysthe validation. Changes of formulation, dimension, or
coefficients require explicit separate model identities.
